\clearpage
\fancyhead[L]{\small GAUSS TO BODE.}
\hypertarget{native-298-lower}{}
\begin{center}
GAUSS TO BODE.
\end{center}

\begin{center}
\begin{tabular}{c}
\hline
Astronomisches Jahrbuch für 1811. Pp. 135--139. Berlin 1808.\\
\hline
\end{tabular}
\end{center}

\begin{flushright}
Göttingen, 24 January 1808.
\end{flushright}

{\rightskip=0pt plus 1em\par\indent The upheavals that first accompanied my change of residence; and the hope I then had of sending you many observations of the comet from here in return for those you shared, a hope that extremely unfavorable weather frustrated, will excuse my long silence. Meanwhile, I send Your --- the few observations of the comet I managed to wrest from the almost constantly overcast sky here.\par}

\begin{center}\fontsize{9}{11}\selectfont\setlength{\tabcolsep}{2pt}
\begin{tabular}{lrrrl}
& M. T. in Göttingen. & R.A. & Decl. & \\
1807. Dec. 17. & $11^{\text{u}}\,24^{\text{m}}\,44^{\text{s}}$ & $305^\circ\;54'\;0''$ & $45^\circ\;4'\;30''$ & \\
29. & $7\;59\;21$ & $318\;57\;41$ & $46\;41\;48$ & \\
& $11\;36\;33$ & $319\;6\;35$ & $46\;43\;11$ & \\
1808. Jan. 4. & $7\;22\;12$ & $325\;11\;19$ & $47\;12\;46$ & \smash{\raisebox{16.5pt}{$\left.\vphantom{\begin{array}{c}0\\0\\0\\0\end{array}}\right\}$ by me.}}\\
& $9\;13\;57$ & $325\;15\;53$ & & \\
& $17\;30\;26$ & $325\;36\;38$ & & by Prof. \textsc{Harding}.\\
\end{tabular}
\end{center}

\par\indent I had already calculated the comet’s parabolic orbit early in November from a few observations made in Bremen and Lilienthal. Since those observations have yet to become known, I pass them on to you here.

\SourceLeaf{299}\noindent Perihelion passage, 18 Sept. $19^{\text{u}}\,16^{\text{m}}\,0^{\text{s}}$ Paris meridian.


\begin{center}\fontsize{10}{12}\selectfont\setlength{\tabcolsep}{3pt}\renewcommand{\arraystretch}{1.15}\begin{tabular}{lr}
Longitude of perihelion .\;.\;.\;.\;.\;.\;.\;.\;.\;.\;.\;.\;.\; & $271^\circ\;0'\;13''3.$ \\
Logarithm of the least distance , .\;.\;.\;.\;.\;.\;.\;.\;.\;.\; & $9.8114927$ \\
Ascending node .\;.\;.\;.\;.\;.\;.\;.\;.\;.\;.\;.\;.\; & $266^\circ\;38'\;31''2.$ \\
Inclination of the orbit .\;.\;.\;.\;.\;.\;.\;.\;.\;.\;.\;.\;.\;.\; & $63\;12\;35.6.$ \\
\end{tabular}\end{center}

\noindent Motion prograde.


\par\indent At the end of November and the beginning of December, these elements differed from the observations by about $1'$, but the error decreased again afterward. If I manage to obtain a few more observations of the comet, I will use them to improve the orbit given above once more, though the corrections will certainly be very small. The observations will hardly settle anything about the orbit’s deviation from a parabola, which cannot be considerable.

\par\indent Here are my new (X.) elements of \emph{Pallas}:


\begin{center}\fontsize{10}{12}\selectfont\setlength{\tabcolsep}{3pt}\renewcommand{\arraystretch}{1.15}\begin{tabular}{lr}
Epoch 1808, Göttingen meridian .\;.\;.\;.\;.\;.\;.\; & $252^\circ\;32'\;28''5$ \\
Daily tropical motion .\;.\;.\;.\;.\;.\;.\;.\;.\;.\;.\;.\; & $770''2143$ \\
Logarithm of the semimajor axis .\;.\;.\;.\;.\;.\;.\;.\;.\; & $0.4423149$ \\
Eccentricity .\;.\;.\;.\;.\;.\;.\;.\;.\;.\;.\;.\;.\;.\;.\;.\;.\; & $0.2450198$ \\
Perihelion 1803 .\;.\;.\;.\;.\;.\;.\;.\;.\;.\;.\;.\;.\;.\;.\; & $121^\circ\;3'\;11''4$ \\
$\AscendingNode\;1803$ .\;.\;.\;.\;.\;.\;.\;.\;.\;.\;.\;.\;.\;.\;.\;.\;.\;.\; & $172\;28\;56.9$ \\
Inclination of the orbit .\;.\;.\;.\;.\;.\;.\;.\;.\;.\;.\;.\;.\;.\; & $34\;37\;41.0$ \\
\end{tabular}\end{center}


\par\indent I recently improved the elements of \emph{Juno} using Mr. \textsc{Bessel}’s observations from last year. Here are the new (VII.) elements:


\begin{center}\fontsize{10}{12}\selectfont\setlength{\tabcolsep}{3pt}\renewcommand{\arraystretch}{1.15}\begin{tabular}{lr}
Epoch, Göttingen meridian, 1805 .\;.\;.\;.\;.\;.\;.\; & $42^\circ\;37'\;3''7$ \\
1806 .\;.\;.\;.\;.\;.\;.\; & $125\;7\;57.8$ \\
1807 .\;.\;.\;.\;.\;.\;.\; & $207\;38\;51.8$ \\
1808 .\;.\;.\;.\;.\;.\;.\; & $290\;23\;19.8$ \\
Mean daily tropical motion .\;.\;.\;.\;.\;.\;.\;.\;.\; & $813''8468$ \\
Tropical period .\;.\;.\;.\;.\;.\;.\;.\;.\;.\;.\; & $1592\;\text{days}\;10\tfrac12\;\text{hr.}$ \\
Perihelion 1805 .\;.\;.\;.\;.\;.\;.\;.\;.\;.\;.\;.\;.\;.\;.\; & $53^\circ\;19'\;0''2$ \\
Ascending node 1805 .\;.\;.\;.\;.\;.\;.\;.\;.\;.\;.\; & $171\;4\;28.2$ \\
Eccentricity .\;.\;.\;.\;.\;.\;.\;.\;.\;.\;.\;.\;.\;.\;.\;.\;.\; & $0.2554996$ \\
Logarithm of the semimajor axis .\;.\;.\;.\;.\;.\;.\;.\;.\; & $0.4263781$ \\
Inclination of the orbit .\;.\;.\;.\;.\;.\;.\;.\;.\;.\;.\;.\;.\;.\; & $13^\circ\;4'\;26.2$ \\
\end{tabular}\end{center}


\par\indent The ephemeris for 1808 calculated from these elements is as follows: [It gives right ascension and declination to the nearest minute for midnight in Göttingen at four-day intervals from 16 April to 28 December.]

\par\indent My work on determining the orbits of celestial bodies is now in press. I hope it can appear by Easter.

